\colorlet{BLACK}{black}
\chapter{\textcolor{black}{Computational Workflow Scripts and Data Structures}}

\section{First Appendix (Tcl Example)}
You can use \texttt{tcblisting} for creating the code snippets. The following example illustrates how one can customize the \texttt{tcblisting} to achieve the \texttt{tcl} script. Similarly, one can use it for other programming language listing, including HDL.

\begin{tcblisting}{listing only,colback=gray!10!white, breakable, boxsep=0pt,top=1mm,bottom=1mm,left=1mm,right=1mm,
listing options={
language= tcl,
basicstyle=\small\ttfamily, 
otherkeywords={create_clock, set_clock_latency, set_input_delay, set_output_delay, set_load, set_max_fanout, set_fanout_load},
keywordstyle=\color{blue}, 
%keywordstyle=[2]{\color{red}},
commentstyle=\color{gray},
backgroundcolor=\color{gray!25},
%morekeywords=[2]{arg,pos},
%moredelim=[is][\color{violet}]{''}{''},
escapechar=!}}
# Since our design has a clock with name clk, 
## specify that name under [get_port ]
create_clock -period 40 -waveform {0 20} [get_ports clk]

# Setting a 'delay' on the clock:
set_clock_latency 0.3 clk

# Setting up constraints on your I/P and O/P pins
set_input_delay 2.0 -clock clk [all_inputs]!\tikz[remember picture]\node[](c1){};!
set_output_delay 1.65 -clock clk [all_outputs]!\tikz[remember picture]\node[](c2){};!

# Set realistic 'loads' on each output pin
set_load 0.1 [all_outputs]

# Set 'maximum' fanin and fan-out for the input and output pins 
set_max_fanout 1 [all_inputs]
set_fanout_load 8 [all_outputs]      
\end{tcblisting}

This appendix also provides pseudocode snippets and data structure examples representing the hybrid Langmuir--Machine Learning computational modelling workflow developed in Python.

\section{Phase 1: Design of Experiments and Simulation}
The following pseudocode illustrates the generation of the Latin Hypercube Sampling (LHS) matrix and the physics-based simulation of fluoride adsorption.

\begin{tcblisting}{listing only,colback=gray!10!white, breakable, boxsep=0pt,top=1mm,bottom=1mm,left=1mm,right=1mm,
listing options={
language=Python,
basicstyle=\small\ttfamily, 
breaklines=true,
keywordstyle=\color{blue}, 
commentstyle=\color{gray},
backgroundcolor=\color{gray!10},
escapechar=!}}
import pandas as pd
from scipy.stats import qmc
import numpy as np

# 1. Generate Latin Hypercube Sampling (LHS) Design
sampler = qmc.LatinHypercube(d=10, seed=42)
sample = sampler.random(n=500)
# Scale to factor ranges (e.g., pH: 3-9, Dose: 0.5-5)
factors = scale_to_ranges(sample) 

# 2. Physics-Based Simulation
def simulate_adsorption(row):
    # Langmuir baseline capacity
    q_langmuir = q_max * KL * row['C0'] / (1 + KL * row['C0'])
    
    # pH effect (Gaussian bell curve centered at 6.5)
    pH_effect = np.exp(-((row['pH'] - 6.5)**2) / 4.5)
    
    # Kinetics, Ion Competition & Fouling penalties
    kinetics = 1 - np.exp(-k * row['Dose'] * row['Time'])
    ion_penalty = calculate_ion_competition(row['Chloride'], 
                                            row['Hardness'], 
                                            row['Carbonate'])
    fouling = 1 - (0.15 * row['NOM'] / 50)
    
    # Final response with 5% Gaussian noise
    q_rem = q_langmuir * pH_effect * kinetics * (1-ion_penalty) * fouling
    return q_rem + np.random.normal(0, 0.05 * q_rem)

df = pd.DataFrame(factors, columns=factor_names)
df['q_removal'] = df.apply(simulate_adsorption, axis=1)
df.to_csv('dataset_simulated_500.csv')
\end{tcblisting}

\section{Phase 2: Langmuir Baseline Model Fitting}
This snippet demonstrates the fitting of the 66-feature polynomial Langmuir baseline model and the calculation of residuals.

\begin{tcblisting}{listing only,colback=gray!10!white, breakable, boxsep=0pt,top=1mm,bottom=1mm,left=1mm,right=1mm,
listing options={
language=Python,
basicstyle=\small\ttfamily, 
breaklines=true,
keywordstyle=\color{blue}, 
commentstyle=\color{gray},
backgroundcolor=\color{gray!10},
escapechar=!}}
from sklearn.preprocessing import PolynomialFeatures, StandardScaler
from sklearn.linear_model import LinearRegression

# Load dataset
df = pd.read_csv('dataset_simulated_500.csv')
X = df[factor_names]
y_actual = df['q_removal']

# Generate 66 features (linear, squared, interactions)
poly = PolynomialFeatures(degree=2, include_bias=False)
X_poly = poly.fit_transform(X)

# Standardize and fit Ordinary Least Squares (OLS)
scaler = StandardScaler()
X_scaled = scaler.fit_transform(X_poly)

model = LinearRegression()
model.fit(X_scaled, y_actual)

# Calculate predictions and residuals
df['q_predicted'] = model.predict(X_scaled)
df['residual'] = df['q_removal'] - df['q_predicted']

df.to_csv('langmuir_predictions.csv')
\end{tcblisting}

\section{Phase 3: Feature Engineering and ML Preparation}
This section illustrates the extraction of engineered features targeted at the non-linear pH residual patterns.

\begin{tcblisting}{listing only,colback=gray!10!white, breakable, boxsep=0pt,top=1mm,bottom=1mm,left=1mm,right=1mm,
listing options={
language=Python,
basicstyle=\small\ttfamily, 
breaklines=true,
keywordstyle=\color{blue}, 
commentstyle=\color{gray},
backgroundcolor=\color{gray!10},
escapechar=!}}
# Feature engineering targeting the Langmuir residual patterns
def engineer_features(df):
    # pH Deviation Features
    df['pH_abs_dev'] = abs(df['pH'] - 6.5)
    df['pH_dev_sq'] = (df['pH'] - 6.5)**2
    df['pH_optimal'] = df['pH'].apply(lambda x: 1 if 6 <= x <= 7 else 0)
    
    # Ratios and Interactions
    df['C0_Dose_ratio'] = df['C0'] / df['Dose']
    df['total_ions'] = df['Chloride'] + (df['Hardness']/10) + \
                       df['Carbonate']
    df['carbonate_at_pH'] = df['Carbonate'] * (df['pH'] - 6.5)
    
    # Logarithmic transformations
    df['log_Time'] = np.log(df['Time'])
    
    return df

ml_data = engineer_features(df)
# Split into training (80%) and test (20%) datasets
ml_data.to_csv('ml_training_data.csv')
\end{tcblisting}

\section{Data File Structures (CSV Snippets)}

\textbf{1. dataset\_simulated\_500.csv (Sample Data)}
\begin{verbatim}
Run, pH, C0_mgL, Time_min, Dose_gL, Temp_C, ..., q_removal
1,  7.00,  70.0,      30.0,    1.00,  30.0, ..., 4.12
2,  2.00,  70.0,     180.0,    0.10,  25.0, ..., 1.45
3,  4.50,   2.0,     105.0,    0.55,  35.0, ..., 3.56
... (500 rows)
\end{verbatim}

\textbf{2. langmuir\_predictions.csv (Model Output)}
\begin{verbatim}
Run, pH, C0_mgL, ..., q_removal, q_predicted, residual
1,  7.00,  70.0, ...,      4.12,        3.47,     0.65
2,  2.00,  70.0, ...,      1.45,        1.97,    -0.52
... (500 rows)
\end{verbatim}